\pg{434} exige necessariamente uma causa proporcionada. A causa \emph{primeira} é \textsc{Deus}. Logo existe \textsc{Deus}\footnote{Os \emph{materialistas}, para destruirem a força d'este argumento, recorrem a uma serie \emph{infinita} de causas \emph{efficientes}, das quaes uma produza outra, sem sermos obrigados a admittir uma \emph{primeira} causa \emph{efficiente}.\par O recurso é vão. Porquanto, essa serie infinita, embora fosse \emph{possivel}. estaria disposta ou em linha \emph{recta}, ou em linha \emph{circular}. Ambas as hypotheses são absurdas. --- E' absurda a \emph{primeira;} porque, se estivesse disposta em linha \emph{recta}, de modo que a causa inferior fosse produzida pela superior, essa serie não teria uma causa \emph{primeira} (porque o infinito não tem principio). Não tendo uma causa \emph{primeira}, não teria a \emph{segunda}, nem a \emph{terceira}, nem causa alguma; visto que numa serie de causas, dispostas em linha \emph{recta}, a inferior não póde existir sem a superior. Logo deve admittir-se, fóra d'esta serie, uma causa primeira. --- E' absurda a segunda hypothese. Na verdade, se a serie \emph{infinita} das causas \emph{efficientes} estivesse disposta em linha \emph{circular}, um e o mesmo ente seria causa e effeito de outro ente. Na verdade, supponhamos uma linha \emph{circular}, formada por a, b, c, d... E' evidente que, sendo a causa de b, b de c, c de d, d de a, d seria causa e effeito (ainda que remoto) de a; o que é absurdo.\par Os adversarios dizem que a \emph{primeira} causa \emph{efficiente} é a \emph{materia}.\par Esta asserção tambem é absurda. Por quanto, como provámos, a materia não é nem póde ser eterna, mas foi produzida, e por isso não póde ser a causa \emph{primeira} de todas as coisas.\par Note-se a differença entre este e o precedente argumento. Aqui, \textsc{Deus} considera-se como a causa \emph{efficiente} de tudo o que existe; além, \textsc{Deus} considerou-se tambem como causa \emph{final} de todas as coisas, emquanto tudo tende para \textsc{Deus}, como Ente infinitamente intelligivel e amavel.}.

c) Os entes visiveis são todos \emph{contingentes}, isto é, existem mas podiam deixar de existir. Se são \emph{contingentes}, foram produzidos por outro ente; porque o \emph{contingente}, podendo existir e deixar de existir, não contém em si a razão sufficiente da propria existencia. Este outro ente, que produziu o \emph{contingente}, ou é tambem \emph{contingente}, ou é \emph{necessario}. Se é \emph{necessario}, é \textsc{Deus}. Se é \emph{contingente}, foi produzido por outro, e este por outro, até se encontrar inevitavelmente o ente \emph{necessario} ; porque, se não existisse o ente \emph{necessario}, não poderia existir nenhum ente \emph{contingente}, pois repugna uma serie de entes \emph{dependentes} sem um \emph{primeiro}, de que todos dependam. Logo existe \textsc{Deus}\footnote{Em poucas palavras: O mundo é \emph{contingente}. Ora o ente \emph{contingente} depende do ente \emph{necessario}. Logo existe um ente \emph{necessario}. O ente necessario é \textsc{Deus}. Logo existe \textsc{Deus}.\par \emph{Vacherot} affirma muitas vezes (sem nunca o demonstrar) que o mundo,\pgn{435}se é \emph{contingente} na \emph{apparencia}, é \emph{necessario} e \emph{absoluto} na \emph{realidade}. --- Esta affirmação é absurda. O mundo, como demonstrámos na \emph{Cosmologia geral}, é \emph{essencialmente contingente} no seu conjuncto e em cada uma das suas partes.\par \emph{Kant} diz que a prova, deduzida da idea do ser \emph{necessario}, não tem valor; porque, como a prova \emph{ontologica}, passa da ordem \emph{ideal} para a \emph{real}, apoiando-se num principio meramente \emph{subjectivo}, que é o de \emph{causalidade}.\par Rejeitamos com \emph{Kant} a prova \emph{ontologica} ácerca da existencia de \textsc{Deus}, mas negamos que a prova \emph{cosmologica} se reduza á \emph{ontologica}. Na verdade, a prova \emph{cosmologica} não consiste, como pensa \emph{Kant}, em deduzir da \emph{idea} de um ser \emph{necessario} a sua existencia \emph{real}, mas em deduzir da existencia \emph{real} dos seres \emph{contingentes} a existencia \emph{real} do ser \emph{necessario}, que é \textsc{Deus}. --- Além d'isso, negamos que o principio de \emph{causalidade} seja meramente \emph{subjectivo}; elle é tambem, e primeiramente, \emph{objectivo}, porque não é lei do nosso espirito, senão porque é lei das coisas.}.
